\chapter{Note sur les entretiens avec M. Girardin}
Cette annexe traite les différents entretien réaliser avec M.Girardin. Différentes notes ont été prise afin d'avoir une trace des conseils données par M. Girardin.


\subsection*{18 février 2026}

\begin{itemize}
    \item Lecture et compréhension du travail précédement réaliser

    \item Se pencher sur le code pour trouver le problème de l'évitation

    \item Réaliser le diagramme d'état

    \item Regarder ce qui a été fait pour l'interface web local

    \item Réaliser une régulation en PID
\end{itemize}

\subsection*{25 février 2026}

\begin{itemize}
    \item Réaliser un organigramme pour les machines d'états des inducteurs.

    \item Réaliser le Gantt pour avoir une vision globale du projet.
\end{itemize}

\subsection*{2 mars 2026}

\begin{itemize}
    \item Vérifier le temps total pris par l'intéruption de la conversions, de la régulation, du PWM et du bouton. Pour cela utiliser l'activation ou l'intéruption d'une LED et de l'oscillo. Faire la démarche avec et sans la communication sans fil.

    \item Vérifier le temps de chaque procesus (Conversion, régulation, PWM,...).

    \item Comprendre comment la communication sans fil fonctionne.

    \item Vérifier le temps de l'envoie et la réception de la communication sans fil au moyen de la LED.

    \item Changer les grands titres du Gantt pour mieux définir les objéctifs à réaliser.
    \item Choisir un code de couleur pour savoir l'avancé du projet p. ex. rouge pour dire que je suis en retard, bleau pour dire que je suis dans les temps, gris pour le temps que je pense faire.
\end{itemize}

\subsection*{16 mars 2026}
Pour les vérifications de temps :
\begin{itemize}
    \item Vérifier l'interruption totale.
    \item Vérifier par bloc (ADC, régul et PWM)
    \item Faire le meme travail mais uniquement pour la régulation de l'avant.
\end{itemize}

\subsection*{17 mars 2026}
Le décollage à 4 inducteurs est censé se dérouler de la manière suivante :

\begin{itemize}
    \item Etat 1 : On monte le courant des 4 inducteurs à 3 A (pour ne pas partir de zéro).
    \item Etat 2 : On augmente le courant des inducteurs 1 et 2 avec une rampe (régulation de courant seule).
    \item Etat 3 : Lorsque l'inducteur 1 ou l'inducteur 2 commence à bouger, on les passe les deux en régulation de position (lévitation de l'avant du véhicule).
    \item Etat 4 : On augmente le courant des inducteurs 3 et 4 avec une rampe (régulation de courant seule).
    \item Etat 5 : Lorsque l'inducteur 3 ou l'inducteur 4 commence à bouger, on les passe les deux en régulation de position (lévitation de l'arrière du véhicule, l'avant lévite déjà).
\end{itemize}


